%% kapa: exam=1 round=29 year=2018 period=2 form=A subject=accounting subject_ko=회계학 seq=39 number=79
%% source: Q-Net 감정평가사 기출문제 (cst00309) artlSeq=5202808 file=extracted/1008/5202808/2018년 제29회 감정평가사 1차시험 기출문제/2018년 제29회 감정평가사 1차시험 기출문제/2교시 A형.pdf
%% conversion: kordoc 4.24.51 -> tools/gen1.py (auto); review=proofread-sonnet; flags=none
\begin{question}{79}
(주)감평은 종합원가계산제도를 채택하고 단일제품을 생산하고 있다. 재료는 공정이 시작되는 시점에서 전량 투입되며, 가공(전환)원가는 공정 전체에 걸쳐 균등하게 발생한다. 가중평균법과 선입선출법에 의한 가공(전환)원가의 완성품환산량은 각각 108,000단위와 87,000단위이다. 기초재공품의 수량이 70,000단위라면 기초재공품 가공(전환)원가의 완성도는?
\begin{choices}
  \item 10\%
  \item 15\%
  \item 20\%
  \item 25\%
  \item 30\%
\end{choices}
\end{question}
\answer{5}
